\clearpage
\section{Gantt Şeması}
\label{gantt}

Çalışma planına ait iş paketleri (tanımları) aşağıdaki gibi tanımlanmalı ve Gantt şeması  kullanılarak gösterilmelidir. En az 4 adet iş paketi tanımlanmalıdır. İş paketlerinde proje grubu elemanlarının hangi görevi yapacakları belirtilmelidir.
\begin{enumerate}
\item İP.1. İş Paketi Adı: Açıklaması
\item İP.2. İş Paketi Adı: Açıklaması
\end{enumerate}


\begin{table*}[!h]
\centering
    \renewcommand{\arraystretch}{1.5}
        \begin{tabular}{| c | c | c | c | c | c | c | c | c | }
              \hline
              \multirow{2}{*}{\parbox{0in}\textbf{{Aylar}}}
              & & & & & & & & \\ \cline{2-9}
              & 11 & 12 & 1 & 2 & 3 & 4 & 5 & 6 \\ \hline

              \multirow{5}{*}{\parbox{0in}\textbf{{İş Tanımı}}}
              						& \multicolumn{3}{c}{\textbf{İP.1}} & & & & & \\ \cline{2-9}
              						& & \multicolumn{4}{|c|}{\textbf{İP.1}} & & & \\ \cline{2-9}
              						& & & & & & & & \\ \cline{2-9}
              						& & & & & & & & \\ \cline{2-9} \hline
    \end{tabular}%}
    \newline
        \caption{Gantt Şeması}
        \label{table:gantt}
    \renewcommand{\arraystretch}{1} 
\end{table*}

\clearpage